% ============================================================
% SNIPPET PARA CAPÍTULO 4 — Análisis de sensibilidad de ventana
% Insertar en subsección "Síntesis de ajustes de rediseño" o
% dentro de "Validación estadística y resiliencia operativa"
% ============================================================

\subsection{Análisis de sensibilidad de la ventana de correlación}
\label{subsec:sensibilidad_ventana}

Antes de fijar los parámetros del contraste principal, se evaluó el impacto de
la ventana de correlación temporal sobre la eficiencia de extracción $\eta$.
El motor utiliza una ventana de 250\,ms por defecto para emparejar eventos
seriales con registros de red; se analizó si ampliarla a 300\,ms alteraba
significativamente los resultados en las sesiones post-intervención.

\begin{table}[H]
\centering
\caption{Efecto de la ventana de correlación sobre $\eta$ en sesiones
         post-intervención (11--26 de mayo de 2026, $n=22$ sesiones).}
\label{tab:sensibilidad_ventana}
\footnotesize
\renewcommand{\arraystretch}{1.22}
\begin{tabularx}{\linewidth}{|>{\raggedright\arraybackslash}p{3.8cm}|>{\centering\arraybackslash}p{2.4cm}|>{\centering\arraybackslash}p{2.4cm}|>{\centering\arraybackslash}p{2.4cm}|>{\centering\arraybackslash}X|}
\hline
\rowcolor{gray!20}
\TableHeaderFirst{Parámetro} & \TableHeader{Ventana 250\,ms} & \TableHeader{Ventana 300\,ms} & \TableHeader{Diferencia} & \TableHeader{Interpretación} \\ \hline

$\bar{\eta}$ media &
20{,}80\,\% &
25{,}96\,\% &
$+5{,}16\,\text{pp}$ &
Ganancia marginal al ampliar la ventana. \\ \hline

Sesiones con mejora &
--- &
--- &
12 de 22 (54{,}5\,\%) &
Menos de la mitad de sesiones mejoran. \\ \hline

Sesiones sin cambio &
--- &
--- &
10 de 22 (45{,}5\,\%) &
Eventos bien alineados o fuera de ambas ventanas. \\ \hline

Decisión metodológica &
\multicolumn{4}{l|}{Se conserva la ventana de 250\,ms como parámetro conservador; la ganancia de 300\,ms no justifica} \\
& \multicolumn{4}{l|}{el riesgo de emparejamientos espurios por coincidencia temporal.} \\ \hline

\end{tabularx}
\end{table}

La elección de 250\,ms se mantiene por dos razones: (i)~la mayoría de los
emparejamientos correctos ocurren dentro de esa ventana, como lo evidencian las
sesiones de validación presencial, y (ii)~ampliar la ventana puede introducir
coincidencias espurias entre eventos de tandas distintas cuando la actividad de
ordeño es intensa. En consecuencia, todos los cómputos de $\eta$ reportados en
el Capítulo~\ref{chapter:chapter06} usan la ventana de 250\,ms.
